\chapter{Background and Related Work}
\label{chap:background}

This chapter develops the concepts needed to compare MI between two
categorical populations. It begins with entropy, pointwise MI, and population
MI. It then reviews plug-in estimation, first-order inference, the
Welch--Satterthwaite approximation, and prior comparisons of information
quantities. The final section identifies the gap addressed by the expanded
method.

\section{Information-theoretic foundations}

\subsection{Entropy and information content}

For a discrete random variable \(X\) with probability mass function
\(p_X(x)\), the information content of outcome \(x\) is
\begin{equation}
  h_X(x)=-\log p_X(x).
\end{equation}
Rare outcomes have larger information content because they are less expected.
Using natural logarithms gives information in nats. The Shannon entropy is the
expected information content,
\begin{equation}
  H(X)=-\sum_x p_X(x)\log p_X(x).
  \label{eq:entropy}
\end{equation}
Entropy summarises the uncertainty of one random variable before its value is
observed \citep{shannon1948}.

For a pair \((X,Y)\), the joint entropy \(H(X,Y)\) is defined from the joint
probabilities in the same way. The conditional entropy
\(
H(X\mid Y)=H(X,Y)-H(Y)
\)
measures the uncertainty remaining in \(X\) after observing \(Y\).

\subsection{Pointwise and average mutual information}

Let \(P=(p_{ij})\) be the joint distribution of categorical variables
\(X\in\{1,\ldots,r\}\) and \(Y\in\{1,\ldots,c\}\). Write
\begin{equation}
  p_{i+}=\sum_{j=1}^{c}p_{ij},
  \qquad
  p_{+j}=\sum_{i=1}^{r}p_{ij}
\end{equation}
for the row and column margins. The pointwise mutual information associated
with cell \((i,j)\) is
\begin{equation}
  \pmi_P(i,j)
  =\log\left(\frac{p_{ij}}{p_{i+}p_{+j}}\right).
  \label{eq:background-pmi}
\end{equation}
The denominator is the cell probability under independence with the same
margins. A positive value means that the cell occurs more often than
independence predicts; a negative value means that it occurs less often.

Mutual information is the probability-weighted average of these pointwise
values:
\begin{equation}
  I(P)
  =\sum_{i=1}^{r}\sum_{j=1}^{c}p_{ij}\pmi_P(i,j).
  \label{eq:background-mi}
\end{equation}
Equivalently,
\begin{equation}
  I(X;Y)=H(X)+H(Y)-H(X,Y).
  \label{eq:mi-entropy-form}
\end{equation}
The quantity is nonnegative and equals zero exactly when
\(p_{ij}=p_{i+}p_{+j}\) for every cell with positive support.

Equation~\eqref{eq:background-mi} is also the Kullback--Leibler divergence
between the joint distribution and the product of its margins. MI therefore
measures the overall departure from independence rather than one particular
shape of association.

This scalar summary deliberately discards much of the information in the
joint distribution. Two populations can have the same MI while assigning
probability to different cells, having different margins, or exhibiting
different association patterns. Equality of MI is therefore weaker than
equality of distributions. This distinction is central to the present test:
the null fixes one scalar functional while leaving the remaining features of
each population unrestricted.

The logarithm base determines the unit but not the inferential problem. If
base two is used, every pointwise value and MI value is divided by
\(\log 2\). The standard error is divided by the same constant, so a
studentized statistic is unchanged. This thesis uses natural logarithms and
reports MI in nats.

\section{Estimating discrete mutual information}

Suppose a sample of size \(n\) produces counts \(N_{ij}\). The empirical cell
probabilities are
\begin{equation}
  \widehat p_{ij}=\frac{N_{ij}}{n}.
\end{equation}
Substitution into Equation~\eqref{eq:background-mi} gives the plug-in estimator
\begin{equation}
  \Ihat(P)
  =\sum_{i,j:\widehat p_{ij}>0}
  \widehat p_{ij}
  \log\left(
  \frac{\widehat p_{ij}}
       {\widehat p_{i+}\widehat p_{+j}}
  \right).
  \label{eq:plugin-mi-background}
\end{equation}
The zero-cell convention follows from
\(\lim_{x\downarrow0}x\log x=0\). Empty cells contribute no direct term, but
they still affect the empirical margins and the finite-sample behaviour of the
estimator.

Under multinomial sampling, the entire table is random subject to its fixed
total. Cell counts are dependent because increasing one count leaves fewer
observations for the others. The plug-in estimator remains consistent for a
fixed alphabet, but its finite-sample distribution depends on the complete
probability table. Equal alphabet sizes and sample sizes do not imply equal
sampling behaviour when one population has much rarer cells.

The plug-in estimator is upward biased. Finite sampling makes an empirical
joint table appear more structured than its population distribution, which
increases the estimated departure from independence. For a fixed positive
\(r\times c\) support, the leading bias is
\begin{equation}
  \Bias\{\Ihat(P)\}
  =\frac{(r-1)(c-1)}{2n}+O(n^{-2}).
  \label{eq:mi-leading-bias}
\end{equation}
This result follows from the corresponding leading entropy biases and is
closely related to Miller's correction \citep{miller1955,moddemeijer1989}.
More complete analyses show that sparse sampling and unknown support can
produce substantial higher-order error beyond Equation
\eqref{eq:mi-leading-bias} \citep{paninski2003}.

The leading term can be seen from the entropy identity in Equation
\eqref{eq:mi-entropy-form}. For a discrete variable with \(m\) positive
categories, the plug-in entropy has leading bias \(-(m-1)/(2n)\). Applying
this result to \(H(X)+H(Y)-H(X,Y)\) gives
\begin{align}
  \Bias\{\Ihat(P)\}
  &\approx
  -\frac{r-1}{2n}
  -\frac{c-1}{2n}
  +\frac{rc-1}{2n}\\
  &=\frac{(r-1)(c-1)}{2n}.
  \label{eq:mi-bias-from-entropies}
\end{align}
The correction is positive because estimating the larger joint distribution
introduces more entropy bias than estimating the two margins.

Bias matters directly when comparing two populations. If
\(I(P)=I(Q)\) but \(n_P\ne n_Q\), the uncorrected difference has leading bias
\begin{equation}
  \frac{(r-1)(c-1)}{2}
  \left(\frac{1}{n_P}-\frac{1}{n_Q}\right).
\end{equation}
Removing the leading term is therefore part of the effect estimate rather than
an optional refinement of its reference distribution.

\section{First-order inference for nonlinear functionals}

MI is a nonlinear function of the cell probabilities. When all population
cells are positive and the empirical table is close to its population value,
a first-order Taylor expansion replaces this nonlinear function locally by a
linear combination of the cell-probability errors. The resulting expansion is
often expressed through an influence function \citep{kandasamy2015}.

For MI, the observation-level influence function is centred pointwise MI:
\begin{equation}
  \psi_P(i,j)=\pmi_P(i,j)-I(P).
  \label{eq:mi-if-background}
\end{equation}
It has population mean zero. Its variance is
\begin{equation}
  V(P)
  =\Var_P\{\psi_P(X,Y)\}
  =\sum_{i,j}p_{ij}\{\pmi_P(i,j)-I(P)\}^2.
  \label{eq:mi-variance-background}
\end{equation}
The plug-in MI error is, to first order, the sample mean of these influence
values. Consequently,
\begin{equation}
  \Var\{\Ihat(P)\}\approx\frac{V(P)}{n}.
  \label{eq:first-order-mi-var-background}
\end{equation}

This representation has two roles. First, it makes the central limit theorem
available because the leading estimation error is an average of independent,
mean-zero contributions. Second, it identifies the observation-level
quantity whose variance determines the standard error. The raw categories
themselves do not have a numerical variance relevant to MI; their effects on
the MI functional are represented by centred pointwise MI.

The expansion is local. It describes sampling fluctuations near a fixed
positive-support population and does not assert that the plug-in estimator is
normally distributed for every finite table. Its accuracy depends on the
sample size relative to the probability structure, not only on the total
number of observations.

The central limit theorem then gives a normal approximation when
\(V(P)>0\). At exact independence, every pointwise-MI value is zero and
\(V(P)=0\). The first-order term disappears, so a second-order expansion is
required. This is why the chi-squared limit for testing independence and the
normal limit for comparing positive MI values are not competing descriptions
of the same regular case.

\section{Wald and Welch--Satterthwaite inference}

\subsection{Normal Wald inference}

A Wald test estimates a parameter difference, divides it by an estimated
standard error, and compares the result with its asymptotic normal reference.
For two independent MI estimates, the first-order variance contributions add:
\begin{equation}
  \Var\{\Ihat(P)-\Ihat(Q)\}
  \approx
  \frac{V(P)}{n_P}+
  \frac{V(Q)}{n_Q}.
  \label{eq:wald-variance-background}
\end{equation}
Replacing the population quantities by estimates produces the normal Wald
test developed formally in Chapter~\ref{chap:baselines}.

Wald inference is asymptotically valid under regularity conditions, but its
finite-sample calibration depends on both the quality of the linear expansion
and the stability of the estimated standard error. A normal reference does not
separately represent the uncertainty of the two estimated variance
contributions.

\subsection{The Welch--Satterthwaite principle}

Welch's test compares two means when their population variances may differ
\citep{welch1947}. The denominator contains two independently estimated
variance contributions. Satterthwaite's approximation represents their sum by
a scaled chi-squared variable with matching first two moments
\citep{satterthwaite1946}. The resulting effective degrees of freedom are
\begin{equation}
  \nu
  =
  \frac{\left(\sum_i k_i s_i^2\right)^2}
       {\sum_i (k_i s_i^2)^2/\nu_i},
  \label{eq:general-ws-background}
\end{equation}
where \(s_i^2\) are variance estimates, \(k_i\) are positive weights, and
\(\nu_i\) are their component degrees of freedom.

The final statistic uses a Student distribution rather than a normal
distribution. Small effective degrees of freedom produce heavier tails,
reflecting greater uncertainty in the estimated denominator. As the component
variance estimates become stable, the effective degrees of freedom increase
and the Student reference approaches the normal reference.

For a single scalar restriction, normal and chi-squared Wald notation are
equivalent. If \(T\) is asymptotically standard normal, then \(T^2\) is
asymptotically \(\chi^2_1\). The signed form is retained here because the
parameter is the scalar difference \(I(P)-I(Q)\), its sign is meaningful, and
the Welch refinement replaces the normal reference for \(T\) by a Student
reference. The \((r-1)(c-1)\)-degree chi-squared law used for one-table
independence is a different second-order result.

\section{Prior comparisons of entropy and mutual information}

Hutcheson applied a Welch--Satterthwaite architecture to differences between
Shannon diversity values \citep{hutcheson1970}. The method combines an entropy
difference, analytic variance estimates, and a Student reference. It is the
closest direct predecessor to the present study because entropy and MI are
both nonlinear functionals of categorical probabilities.

The analogy is useful but incomplete. Entropy depends directly on one
probability vector. Discrete MI depends on the joint probabilities and on the
two margins calculated from them. Moving probability into one cell therefore
changes not only that cell's pointwise MI but also pointwise values throughout
its row and column. An MI-specific variance calculation must retain these
indirect effects.

The comparison of MI-related quantities is also established. Mora and
Ruiz-Castillo studied the estimation and comparison of an MI-based segregation
index \citep{mora2009,mora2011}. Their work places MI comparison within a
broader literature on nonlinear inequality and segregation measures. Related
research combines analytic bias adjustments, asymptotic tests, and bootstrap
refinements for such indices \citep{allen2015}.

In neuroscience, MI estimates have been compared using jackknife uncertainty
estimates \citep{hart2010}. Other applications apply ordinary mean tests to a
collection of replicated MI estimates rather than deriving uncertainty from a
single multinomial table in each population. These uses demonstrate the
practical demand for differential MI inference, but they do not provide the
variance-functional degrees of freedom developed here.

Examples of the replicated-estimate approach occur across several fields.
MI values from randomised astronomical samples have been compared with a
Student test \citep{sarkar2020}, and per-neuron information measurements have
been analysed with Welch comparisons \citep{prince2021}. Partitioned
continuous-data MI estimates have also been compared using unequal-variance
mean tests \citep{martin2026}. In each case, replication supplies a sample of
MI values. The present problem instead begins with one multinomial table from
each population and must derive uncertainty from the cell counts themselves.

Information-theoretic statistics and mean-based tests have also been used as
parallel detectors rather than as parts of one inferential construction
\citep{mather2013}. These applications show that the combination of MI and
Welch terminology is not a novelty claim. The narrower question is whether the
variance-component logic can be derived for the plug-in MI functional itself.

Hutcheson's result also defines an important novelty boundary. Substituting MI
for entropy in a Welch-like template is not by itself a new statistical
principle. MI introduces additional structure because its pointwise values
depend on both joint and marginal probabilities. The sampling uncertainty of
the resulting MI variance estimate must account for those changing margins.

\section{Resampling and the weak null}

Bootstrap methods can approximate the sampling distribution of nonlinear
indices by repeatedly sampling from fitted empirical populations. They are
flexible but computationally more expensive than an analytic table scan.
Their accuracy also depends on how well the empirical distribution represents
rare cells and unknown support.

The bootstrap and permutation address different sources of uncertainty. A
multinomial bootstrap samples each fitted table separately and approximates
the estimator's repeated-sampling distribution. A label permutation pools the
observations and reallocates them between groups. The latter construction is
exact only when the pooled observations are exchangeable under the null.
Bootstrap refinements have improved inference for nonlinear segregation and
entropy measures, which makes them valuable future comparators
\citep{allen2015,palma2022}. They do not remove the need to state which null is
being approximated.

Permutation requires a separate distinction. Pooling two samples and
permuting their group labels is exact under the strong null \(P=Q\), because
the observations are exchangeable. Under the weak null \(I(P)=I(Q)\), the two
complete distributions may differ, so raw label permutation generally targets
the wrong null. Studentized permutation can recover asymptotic validity for
some weak-null parameter comparisons \citep{chung2013}, but it is not the
standard one-table MI shuffle used to test independence.

The primary thesis comparison therefore uses two transparent analytic
baselines that share the same effect estimate and standard error as the
proposed method. Resampling remains relevant for future confirmatory work, but
it is not treated as an exact finite-sample oracle for the unrestricted weak
null.

\section{Literature synthesis and methodological gap}

The literature establishes all of the broad ingredients used in this thesis:
plug-in MI estimation and bias correction, first-order influence-function
variance, Wald comparison of smooth functionals, and
Welch--Satterthwaite moment matching. It also establishes prior comparisons of
entropy and MI-related indices.

No single ingredient is presented as new. Shannon established MI; Miller and
later authors developed finite-sample bias analysis; influence-function and
delta-method theory provide first-order variance; and Welch, Satterthwaite,
and Hutcheson provide the variance-component architecture. Prior MI comparison
work establishes the scientific problem. The thesis contribution is the
particular MI variance-functional derivation needed to connect these pieces,
followed by a controlled evaluation under distinct equal-MI populations.

The remaining gap is narrower. The ordinary Welch substitution assigns a
classical sample-variance degree of freedom, such as \(n-1\), to each variance
component. The estimated MI variance is not an ordinary sample variance of
fixed observations. It is a nonlinear function of the same empirical joint
table used to calculate MI. A change in one observation modifies the cell
probabilities, margins, pointwise-MI values, and their weighted variance.

The expanded method addresses this gap by differentiating the complete MI
variance functional. The resulting observation-level sensitivity determines
the sampling variance of the variance estimate and, through Satterthwaite
moment matching, an MI-specific component degree of freedom. Chapters
\ref{chap:baselines} and \ref{chap:expanded} develop this construction.
